% ============================================================================
%  SECTION 2 — Mathematical details / finite-strain fundamentals
% ============================================================================
\begin{frame}[t]{}
  \vfill\begin{center}{\Huge\color{YaleBlue} 2.\ \ Finite-strain fundamentals}\\[0.5em]
  {\color{BoxGray} just enough continuum mechanics to read the four theories}\end{center}\vfill
\end{frame}

\begin{frame}[t]{Deformation, stretch, and (in)compressibility}
  A body deforms from a \textbf{reference} to a \textbf{current} configuration; the
  deformation gradient $\bm{F}$ maps reference line elements to current ones.
  Everything here collapses to one dominant direction (bar axis, or artery
  circumference), so we track a single scalar \textbf{stretch}
  \giveneq{\lambda = \frac{\text{current length}}{\text{reference length}},\qquad
    \lambda=1\ \text{undeformed},\quad \lambda>1\ \text{extension}.}
  The tissue is \textbf{elastically incompressible} --- only \emph{growth} changes
  volume --- so the through-thickness and transverse stretches follow from
  $\lambda$ alone. One number suffices.
  \vspace{0.4em}

  We report the \textbf{Cauchy (true) stress} $\sigma$ (force per unit
  \emph{current} area): the physically meaningful measure for a growing body,
  because the load-bearing area itself changes as the tissue remodels.
  \srccite{finite-strain background adapted from YaleCBL-Teaching/BENG\_5570}
\end{frame}

\begin{frame}[t]{Wall stress from the loading --- the Laplace law}
  For a pressurised thin tube (the artery), the circumferential wall stress is
  \slboxeq{2.6em}{eq:laplace}{\sigma_\theta = \frac{P\,r}{h}}
  with pressure $P$, current inner radius $r$, wall thickness $h$. The bar instead
  carries a stress set by an \emph{axial} load. The consequence, drawn out on the
  right: the ``required stress'' the loading demands grows
  \textbf{linearly} in $\lambda$ for the bar, but \textbf{quadratically} for the
  artery.
  \vspace{0.3em}
  \begin{columns}[T]
    \begin{column}{0.46\textwidth}
      \begin{itemize}\small
        \item bar: $\ \sigma_{\text{req}}\propto \lambda^1 / (M/M_0)$
        \item artery: $\ \sigma_{\text{req}}\propto \lambda^2 / (M/M_0)$
      \end{itemize}
      {\color{AccentRed}\small That extra power of $\lambda$ is the whole story of
      arterial instability (§7).}
    \end{column}
    \begin{column}{0.54\textwidth}
      \vspace{-0.6em}\slidefig[\linewidth]{fig01_constitutive.pdf}
    \end{column}
  \end{columns}
  \figcredit{Right panel reproduced with the \texttt{gr} package (this repo)}
\end{frame}

\begin{frame}[t]{Hyperelastic constituents}
  Each constituent stores energy in a strain-energy function; its stress is the
  derivative. We use the two laws that describe arterial tissue.
  \vspace{0.3em}

  \textbf{Elastin --- neo-Hookean} (soft, barely stiffening):
  \giveneq{\sigma_e(\lambda_e) = c_e\!\left(\lambda_e^2 - \tfrac{1}{\lambda_e}\right).}

  \textbf{Collagen \& smooth muscle --- Fung exponential fiber} (crimped, then
  stiffens steeply once recruited):
  \slboxeq{2.8em}{eq:fung}{\sigma(\lambda_e) = c_1\,\lambda_e^2\,(\lambda_e^2-1)\,
    e^{\,c_2(\lambda_e^2-1)^2},\qquad I_4=\lambda_e^2}
  These are exactly the constituent free energies used in the constrained-mixture
  literature. The tangent $\dd\sigma/\dd\lambda_e$ is also needed --- for
  equilibrium solves and for the stability analysis (§7).
  \srccite{constituent laws as in Latorre, Humphrey, Comput Methods Appl Mech Eng (2020)}
\end{frame}

\begin{frame}[t]{The deposition stretch $G^k$ --- the crucial modelling idea}
  Cells deposit new fibers \textbf{already under tension}, at a fixed
  \textbf{deposition stretch} $G^k>1$. So at homeostasis a constituent's elastic
  stretch equals its own $G^k$, and we \emph{define} its homeostatic stress from
  that set-point:
  \slboxeq{2.4em}{eq:setpoint}{\sigma_h^k := \sigma^k(G^k)}
  This single convention pins down \emph{every} set-point in the course from a few
  physical inputs --- and is why all four theories can share one parameter set and
  give apples-to-apples comparisons.
  \srccite{deposition-stretch concept: Humphrey, Rajagopal, Math Models Methods Appl Sci (2002)}
\end{frame}

\begin{frame}[t]{Multiplicative decomposition --- growth vs.\ elasticity}
  The single most important kinematic idea in G\&R: split the deformation into a
  stress-free \textbf{inelastic} part (growth / remodeling) and an \textbf{elastic}
  part --- the only part that carries stress. In 1D,
  \slboxeq{2.4em}{eq:split}{\lambda = \lambda_e \,\lambda_{\text{inel}}}
  \begin{itemize}
    \item \textbf{Kinematic growth} (§3) uses \emph{one} such split with an
    inelastic \emph{growth} stretch.
    \item \textbf{Constrained mixtures} (§4) give \emph{each cohort of each
    constituent} its own inelastic natural configuration, and superpose them.
  \end{itemize}
  \vspace{0.4em}
  The set-point \eqref{eq:setpoint} and the split \eqref{eq:split} are the two
  pieces of machinery reused in every model that follows.
\end{frame}
